\subsection{An arithmetic-specialization check in the elliptic search}

The search for the converse in 18:13 also led to the complete
18-page preprint \cite{BhowmikPathak2024}. We distinguish that
consulted version from its 23-page journal publication, whose
full text was not obtained. The following correction concerns
the statement printed in the consulted preprint only.

\begin{result}{G.2.2}{A specialization counterexample and the coefficient condition that repairs it}{Refuted; corrected statement proved}
The lower bound in Theorem~1.11 of the consulted version of
\cite{BhowmikPathak2024} is false under its printed hypotheses.
An exact counterexample is
\[
 \tau=i,\qquad f=E_4,\qquad g=j,\qquad
 \widetilde f=E_4+\frac{E_2E_6}{E_4}.
 \tag{G.2.2a}
\]

The following specialization statement holds. Let $f$ be a
nonzero meromorphic modular form of weight $k\ne0$, let $g$
be a nonconstant modular function, and write
\[
 \widetilde f=\sum_{r=0}^p f_r E_2^r,
 \qquad \operatorname{weight}(f_r)=w-2r.
\]
All forms and functions are for congruence subgroups and have
algebraic Fourier coefficients. At $\tau\in\mathbb H$, assume
that $f,g,f_0,\ldots,f_p$ are finite, $f(\tau)\ne0$, and
$f_r(\tau)\ne0$ for at least one $r\ge1$. Then, with
$q_\tau=\exp(2\pi i\tau)$,
\[
\begin{split}
 &\trdeg_{\Qb}\Qb
   \bigl(q_\tau,f(\tau),g(\tau),\widetilde f(\tau)\bigr)\\
 &\qquad=\trdeg_{\Qb}\Qb
   \bigl(q_\tau,E_2(\tau),E_4(\tau),E_6(\tau)\bigr)
   \ge3.
\end{split}
\]
Thus the nonvanishing condition is imposed on an actual
positive-depth coefficient at the point of evaluation.
\end{result}
\begin{proof}
The transformation law for $E_6$ at the fixed point of
$\tau\mapsto-1/\tau$ gives
$E_6(i)=i^6E_6(i)=-E_6(i)$, hence $E_6(i)=0$.
Since $E_4^3-E_6^2=1728\Delta$ and $\Delta$ has no zeros
in $\mathbb H$, $E_4(i)\ne0$ and $j(i)=1728$.
The expression in (G.2.2a) has weight four and depth one:
its coefficient of $E_2$ is the nonzero meromorphic
weight-two form $E_6/E_4$. Every Fourier coefficient is
rational. At infinity it begins $2-528q+O(q^2)$, so it
is nonconstant. At $i$ all values in question are finite,
and $f(i)=\widetilde f(i)=E_4(i)\ne0$. Consequently
\[
 \Qb\bigl(e^{-2\pi},f(i),g(i),\widetilde f(i)\bigr)
       =\Qb\bigl(e^{-2\pi},E_4(i)\bigr),
\]
whose transcendence degree is at most two. This verifies
every stated hypothesis of the consulted Theorem~1.11 and
contradicts its bound three. Its positive-depth coefficient
has become zero at the evaluation point. The rational
expansion and specialization identities are checked in
\texttt{check\_modular\_specialization.py}; the transcendence-degree
upper bound follows from the two displayed generators.

We prove the corrected statement. Pass to a common congruence
subgroup on which all the functions are defined. On its
compact modular curve, both $j$ and the nonconstant function
$g$ give finite morphisms to $\mathbf P^1$. The functions
and the curve are defined over $\Qb$; the function-field
description is \cite[Theorem 2.11, pp.10--11]{BhowmikPathak2024}.
At the point determined by $\tau$, finiteness implies that
the value of every regular modular function is algebraic
over $\Qb(g(\tau))$, and also over $\Qb(j(\tau))$.
This statement includes a transcendental value: its
scheme-theoretic image is then the generic point, and the
same finite extension of residue fields applies.

Put $\delta=\Delta(\tau)\ne0$. The modular function
$f^{12}/\Delta^k$ is finite and nonzero at $\tau$, so its
value is algebraic over $\Qb(g(\tau))$. The equality
\[
 f(\tau)^{12}
   =\bigl(f^{12}/\Delta^k\bigr)(\tau)\,\delta^k
\]
and $k\ne0$ make $\delta$ algebraic over
$\Qb(f(\tau),g(\tau))$. The same field therefore contains
$j(\tau)$, $E_4(\tau)$ and $E_6(\tau)$ algebraically, using
\[
 E_4(\tau)^3=j(\tau)\delta,\qquad
 E_6(\tau)^2=(j(\tau)-1728)\delta.
\]
Applying the same ratio construction to each $f_r$ shows
that its finite value is algebraic over this field; zero
values cause no problem. The specialized polynomial
\[
 \sum_{r=0}^p f_r(\tau)X^r-\widetilde f(\tau)
\]
is nonconstant by the stated coefficient condition, and
vanishes at $X=E_2(\tau)$. Hence $E_2(\tau)$ is algebraic
over $\Qb(f(\tau),g(\tau),\widetilde f(\tau))$.

Conversely the finite map given by $j$, followed by the same
twelfth-power ratio identities, makes $f(\tau),g(\tau)$
and all $f_r(\tau)$ algebraic over
$\Qb(E_4(\tau),E_6(\tau))$. The polynomial expression then
makes $\widetilde f(\tau)$ algebraic over the field obtained
by adjoining $E_2(\tau)$. These two algebraic extensions
prove the displayed equality after adjoining $q_\tau$.
Nesterenko's theorem \cite[Theorem 1]{EllNesterenko1996},
in the precise form used in A.4,
gives the lower bound three.
\end{proof}

This check does not replace $E_2$ by $E_2^*$ in that
theorem. The exact change of value is
$E_2^*(\tau)=E_2(\tau)-3/(\pi\operatorname{Im}\tau)$;
the additional term remains in the Hodge-splitting equation
of A.4.1. The failure of the attempted specialization
therefore leaves the non-CM algebraicity obligation in place.

Kumar's main theorem \cite[p.410, with Lemma 3 on p.412]{Kumar2021}
was also checked against that equation. Its premise is
$\Q$-linear dependence of the ordinary quasiperiods.
For the Hodge class the actual transformed row is
\[
 (\eta_1-h\omega_1,\eta_2-h\omega_2)
   =\frac{\pi}{\operatorname{Area}(\Lambda)}
        (\bar\omega_1,\bar\omega_2).
\]
Its ratio is $\bar\tau$, which is transcendental in the
non-CM algebraic-invariant case by Schneider. Thus this
transformed row satisfies neither the rational-coefficient
nor the algebraic-coefficient dependence premise. The
conclusion on Kumar's p.414 uses $\Qb$, whereas his main
theorem uses $\Q$; the original printed bars were visually
checked, since text extraction loses them.
